% 作品介绍：面向线性PCurve的一类通用样条面到二次曲面智能转化方案
% 第二章：2026-10-16 用户确认稿。此文件为章节片段，待整稿组装。
% 依据 docs/requirements/中船海舟样条转二次曲面.pdf 第1页“命题目标”。
% 以下为命题要求，实际实现范围与达成情况分别在第三章、第六章说明。
% 整稿需加载 array、tabularx 宏包。
\section{问题重述}

本项目面向船舶设计中的样条曲面几何处理需求，研究给定 NURBS 面片的解析曲面可转换性判定与几何表示恢复问题。对于能够在规定容差内采用目标解析曲面表示的输入，应识别其曲面类型，恢复几何参数及有效边界，并验证转换前后的几何一致性；对于不满足转换条件或无法完成可靠判定的输入，应返回明确、可追溯的处理状态。

根据命题要求，该问题需要同时满足以下五类目标。

\begin{table}[htbp]
\centering
\caption{命题目标与评价要求}
\label{tab:problem-objectives}
\begin{tabularx}{\linewidth}{@{}>{\bfseries}p{0.16\linewidth}X@{}}
\hline
目标类别 & 具体要求 \\
\hline
功能目标 & 实现样条面到目标解析曲面的可转换性判定与转换，形成可供几何内核后续处理的曲面表示。 \\
精度目标 & 转换前后曲面的最大几何距离与平均距离误差均不超过容差 $d$，默认 $d=10^{-5}\,\mathrm{mm}$；对应边界曲线之间的 Hausdorff 距离不超过 $10^{-5}\,\mathrm{mm}$；相邻曲面公共边界处的法向夹角不超过 $0.1^\circ$，满足命题规定的 G1 连续性误差要求。 \\
性能目标 & 对控制顶点总数为 50～200 的单个样条面片，转换判定耗时不超过 $10\,\mathrm{ms}$；转换后的投影、求交和布尔运算耗时较原样条表示降低至少 $50\%$。 \\
兼容性目标 & 使用 C++ 实现，基于海星 SGK 1.7 几何内核开发，满足工程集成要求。 \\
稳定性目标 & 连续处理 1,000 个不同复杂度的 NURBS 曲面时，不发生算法崩溃、内存泄漏或结果不一致；对损坏的输入文件、超出有效范围的参数等异常情况，提供清晰、可追溯的错误信息。 \\
\hline
\end{tabularx}
\end{table}

因此，本项目需要解决的不仅是曲面类型识别，还包括容差约束下的几何参数恢复、有限面片边界一致性验证，以及面向连续批量处理的效率与可靠性问题。上述指标共同构成方案设计和实验评价的依据，其中单面片转换判定效率与转换后的下游运算效率应分别测试、独立评价。
